\documentclass{beamer}

\usepackage{beamerthemeSingapore}
\usepackage[brazil]{babel}
\usepackage[T1]{fontenc}
\usepackage[utf8]{inputenc}
\usepackage{graphicx}
\usepackage{url}
\usepackage{xcolor}
\usepackage{booktabs}
\usepackage{ragged2e}
\usepackage{etoolbox}
\usepackage{minted}
\usepackage{tikz}
\usetikzlibrary{positioning, arrows.meta, shapes.geometric}

% Sob LuaTeX, o grupo vazio de "-{}-" não impede a ligadura de dois hífens em
% fonte monoespaçada, e o resultado sai como travessão. O \mbox{} impede.
\newcommand{\ddash}{-\mbox{}-}

\newcommand{\tcb}[1]{\textcolor{blue}{#1}}
\newcommand{\tcr}[1]{\textcolor{red}{#1}}
\newcommand{\tcg}[1]{\textcolor{green!50!black}{#1}}

\newenvironment{noterminal}
  {\begin{exampleblock}{No seu terminal}}
  {\end{exampleblock}}

\newenvironment{nonavegador}
  {\begin{exampleblock}{No navegador}}
  {\end{exampleblock}}

\setminted{
  fontsize=\footnotesize,
  breaklines=true,
  breakanywhere=true,
  frame=lines,
  framesep=2mm,
  baselinestretch=1.05,
  tabsize=2
}

\apptocmd{\frame}{}{\justifying}{}

\setbeamertemplate{navigation symbols}{}

\title{02 - Git e GitLab}
\author{Prof. Alex Kutzke}
\date{Agosto de 2026}

\begin{document}

% ==========================================
\begin{frame}
  \begin{center}
    \large{Setor de Educação Profissional e Tecnológica}\\
    Tecnologia em Análise e Desenvolvimento de Sistemas\\
    \vspace{1cm}
    \small{DS122 - Desenvolvimento de Aplicações Web 1}
  \end{center}
  \titlepage
\end{frame}

% ==========================================
\begin{frame}{Objetivos da aula}
  Ao final desta aula você deve ser capaz de:
  \begin{enumerate}
    \item Explicar que problema o controle de versão resolve;
    \item Descrever as três áreas de um repositório local e dizer em qual delas
          um arquivo está, a partir da saída de \texttt{git status};
    \item Executar o ciclo \texttt{add}, \texttt{commit} e \texttt{push}, e ler
          o que cada comando devolveu;
    \item Escrever mensagens de \textit{commit} que descrevam a alteração feita;
    \item Ler o histórico com \texttt{git log}, \texttt{git show} e
          \texttt{git diff};
    \item Resolver um conflito de mesclagem editando o arquivo marcado pelo Git;
    \item Entregar um exercício da disciplina pelo GitLab, do \textit{fork} ao
          \texttt{push}.
  \end{enumerate}
\end{frame}

% ==========================================
\begin{frame}{Dinâmica da aula}
  Aula de mãos no teclado. Abra agora um terminal e deixe-o visível ao lado dos
  slides.

  \vspace{0.4cm}
  \begin{exampleblock}{No seu terminal}
    Os quadros verdes são o que você digita. Digite junto com o professor e
    espere a explicação antes de seguir para o próximo.
  \end{exampleblock}

  \vspace{0.4cm}
  Entre um comando e outro, paramos para olhar a saída. Ler o que o Git
  respondeu é metade do aprendizado desta aula: quase todo problema com Git na
  disciplina vem de alguém que digitou sem ler a resposta.
\end{frame}

% ==========================================
\begin{frame}{Roteiro}
  \tableofcontents
\end{frame}

% ==========================================
\section{Preparação}
% ==========================================

\begin{frame}[fragile]{O problema}
  A pasta do trabalho da disciplina anterior:

  \begin{minted}{text}
trabalho.odt
trabalho_v2.odt
trabalho_v2_revisado.odt
trabalho_FINAL.odt
trabalho_FINAL_agora_vai.odt
  \end{minted}

  \vspace{0.2cm}
  O arranjo funciona no sentido mais fraco possível: nada foi perdido. Fora
  isso, não responde a nada.

  \begin{itemize}
    \item Qual é a versão mais recente, a \texttt{FINAL} ou a
          \texttt{FINAL\_agora\_vai}?
    \item O que mudou de \texttt{v2} para \texttt{v2\_revisado}?
    \item Em qual arquivo ainda existe o parágrafo apagado por engano na terça?
    \item Se o trabalho fosse em dupla, quem escreveu o quê?
  \end{itemize}
\end{frame}

\begin{frame}{Sistema de controle de versão}
  Guarda o histórico de um conjunto de arquivos. O que ele registra não é o
  conteúdo atual: é a \textbf{sequência de estados} pela qual o conjunto passou,
  com autoria, data e descrição em cada passo.

  \vspace{0.5cm}
  \begin{block}{Git}
    Criado em 2005 por Linus Torvalds para o desenvolvimento do núcleo do Linux.
    Software livre, sistema dominante hoje, e roda inteiramente na sua máquina.
  \end{block}
\end{frame}

\begin{frame}{O que o Git resolve nesta disciplina}
  \textbf{Histórico.} Você quebra o CSS às 22h e não faz ideia do que mexeu.
  Compare o estado atual com o do último \textit{commit} e veja a linha
  alterada.

  \vspace{0.4cm}
  \textbf{Duas máquinas.} Você começa o exercício no laboratório e termina em
  casa. Sem repositório central, a sincronização vira \textit{pen drive} e anexo
  de e-mail, e uma hora as duas cópias divergem.

  \vspace{0.4cm}
  \textbf{Entrega auditável.} O histórico mostra o que você entregou, quando e
  \textbf{como} chegou lá. É por isso que o diário de bordo do trabalho prático
  é avaliado.
\end{frame}

\begin{frame}[fragile]{\texttt{git \ddash version}}
  \begin{noterminal}
    \begin{minted}{bash}
git --version
    \end{minted}
  \end{noterminal}

  \begin{minted}{text}
git version 2.55.0
  \end{minted}

  Se o comando não for encontrado, instale a partir de

  \begin{center}
    \url{https://git-scm.com/downloads}
  \end{center}

  \vspace{0.3cm}
  \begin{alertblock}{Windows}
    O instalador oficial traz o \textbf{Git Bash}, um terminal onde todos os
    comandos desta aula funcionam como escritos. Use-o. No PowerShell, alguns
    exemplos se comportam de outro jeito.
  \end{alertblock}
\end{frame}

\begin{frame}[fragile]{\texttt{git config}}
  \begin{noterminal}
    \begin{minted}{bash}
git config --global user.name "Seu Nome Completo"
git config --global user.email "seuemail@example.com"
git config --global init.defaultBranch main
    \end{minted}
  \end{noterminal}

  O autor de cada \textit{commit} vem daí, e não do serviço onde o repositório
  está hospedado. Sem configuração, o Git inventa algo a partir do usuário do
  sistema, e o histórico fica assinado por \texttt{aluno@lab-a15-pc07}.

  \vspace{0.3cm}
  \begin{alertblock}{Nos computadores do laboratório}
    A máquina é de todo mundo. Execute os dois primeiros comandos
    \textbf{sem} \texttt{\ddash global}, dentro de cada repositório clonado.
  \end{alertblock}
\end{frame}

\begin{frame}[fragile]{Onde cada valor de configuração foi parar}
  \begin{noterminal}
    \begin{minted}{bash}
git config --list --show-origin | grep -E "user|defaultBranch"
    \end{minted}
  \end{noterminal}

  A saída mostra, para cada valor, o arquivo de onde ele veio. Três níveis
  possíveis, do mais geral ao mais específico:

  \begin{itemize}
    \item do sistema, para todos os usuários da máquina;
    \item do usuário, o que \texttt{\ddash global} escreve;
    \item do repositório, dentro do próprio \texttt{.git/config}.
  \end{itemize}

  \vspace{0.2cm}
  O mais específico vence. É por isso que a configuração local do laboratório
  resolve.
\end{frame}

% ==========================================
\section{Ciclo local}
% ==========================================

\begin{frame}[fragile]{\texttt{git init}}
  \begin{noterminal}
    \begin{minted}{bash}
mkdir catalogo
cd catalogo
git init
    \end{minted}
  \end{noterminal}

  \begin{minted}{text}
Initialized empty Git repository in /home/ana/catalogo/.git/
  \end{minted}

  Nenhuma pasta de destino, nenhum servidor, nenhuma conta. Tudo o que o Git fez
  aconteceu dentro de \texttt{catalogo/}.
\end{frame}

\begin{frame}[fragile]{O repositório é a pasta \texttt{.git}}
  \begin{noterminal}
    \begin{minted}{bash}
ls -a
    \end{minted}
  \end{noterminal}

  \begin{minted}{text}
.  ..  .git
  \end{minted}

  Ali dentro mora o histórico inteiro: todas as versões de todos os arquivos,
  todos os autores, todas as datas.

  \begin{itemize}
    \item Copiar a pasta \texttt{catalogo/} copia o repositório inteiro. Apagar
          \texttt{.git} devolve uma pasta comum, sem histórico nenhum.
    \item \textbf{O Git não precisa de rede.} Tudo até a seção do remoto
          funciona com o cabo desconectado.
  \end{itemize}
\end{frame}

\begin{frame}{Repositório dentro de repositório}
  \begin{alertblock}{Não execute \texttt{git init} dentro de um repositório}
    Os dois passam a disputar os mesmos arquivos, e a confusão é difícil de
    desfazer.
  \end{alertblock}

  \vspace{0.5cm}
  Na dúvida sobre onde você está, \texttt{git status} antes de \texttt{git
  init}. Se ele responder alguma coisa em vez de reclamar que não há
  repositório, você já está dentro de um.
\end{frame}

\begin{frame}[fragile]{\texttt{git status}}
  \begin{noterminal}
    \begin{minted}{bash}
echo '<h1>Catálogo de Produtos</h1>' > index.html
git status
    \end{minted}
  \end{noterminal}

  \begin{minted}{text}
On branch main

No commits yet

Untracked files:
  (use "git add <file>..." to include in what will be committed)
	index.html

nothing added to commit but untracked files present
  \end{minted}

  \textbf{Não rastreado} (\textit{untracked}): o Git nunca viu esse arquivo e
  não vai gravá-lo sem ordem explícita.
\end{frame}

\begin{frame}{As três áreas}
  \begin{center}
    \resizebox{\linewidth}{!}{%
    \begin{tikzpicture}[
        box/.style={draw, rounded corners, minimum width=2.9cm,
                    minimum height=1.5cm, align=center, thick},
        every node/.style={font=\scriptsize}
      ]
      \node[box, fill=blue!8]   (w) at (0,0)   {\textbf{Diretório}\\\textbf{de trabalho}\\\textit{working tree}};
      \node[box, fill=orange!12](s) at (4.8,0) {\textbf{Área de}\\\textbf{preparação}\\\textit{staging area}};
      \node[box, fill=green!10] (r) at (9.6,0) {\textbf{Repositório}\\\texttt{.git}};

      \draw[-{Stealth[length=3mm]}, thick, blue]
        (w.east) -- node[above, font=\scriptsize]{\texttt{git add}} (s.west);
      \draw[-{Stealth[length=3mm]}, thick, green!50!black]
        (s.east) -- node[above, font=\scriptsize]{\texttt{git commit}} (r.west);
    \end{tikzpicture}}
  \end{center}

  \begin{itemize}
    \item \textbf{Diretório de trabalho}: a pasta como o editor de texto a vê.
    \item \textbf{Área de preparação}: a lista do que vai entrar no próximo
          \textit{commit}.
    \item \textbf{Repositório}: o histórico já gravado.
  \end{itemize}
\end{frame}

\begin{frame}{Para que serve a área de preparação}
  É a parte que causa mais estranheza no primeiro contato, porque parece um
  passo burocrático a mais.

  \vspace{0.4cm}
  \begin{block}{}
    Ela existe para você escolher \textbf{o que entra em cada commit}.
  \end{block}

  \vspace{0.4cm}
  Você mexeu em cinco arquivos: dois resolvem um problema e três resolvem outro.
  Dá para gravar dois \textit{commits} separados, cada um com sua mensagem, em
  vez de um pacote único e ininteligível.
\end{frame}

\begin{frame}[fragile]{\texttt{git add} e \texttt{git commit}}
  \begin{noterminal}
    \begin{minted}{bash}
git add index.html
git status
git commit -m "Cria página inicial do catálogo"
    \end{minted}
  \end{noterminal}

  Entre um comando e outro, leia: o \texttt{index.html} saiu de
  \texttt{Untracked files} e apareceu sob \texttt{Changes to be committed}.

  \begin{minted}{text}
[main (root-commit) 5d09b83] Cria página inicial do catálogo
 1 file changed, 1 insertion(+)
 create mode 100644 index.html
  \end{minted}
\end{frame}

\begin{frame}{O que um commit guarda}
  Um \textit{commit} é o registro de um \textbf{estado completo} do projeto, e
  não a lista do que mudou. O Git guarda a fotografia de todos os arquivos
  naquele instante.

  \begin{itemize}
    \item o conteúdo de todos os arquivos versionados;
    \item o \textbf{autor}, com nome e e-mail;
    \item a \textbf{data};
    \item a \textbf{mensagem} que você escreveu;
    \item o \textbf{commit anterior}, chamado de \textit{pai};
    \item um \textbf{identificador} de 40 caracteres hexadecimais, o
          \textit{hash}.
  \end{itemize}

  \vspace{0.2cm}
  \footnotesize Internamente o Git evita duplicar o conteúdo de arquivos que não
  mudaram, mas isso é otimização de armazenamento e não altera o modelo.
\end{frame}

\begin{frame}{O identificador é calculado a partir de todo o resto}
  Conteúdo, autor, data, mensagem e pai entram na conta. Alterar qualquer um
  produz um identificador diferente.

  \vspace{0.4cm}
  Daí vêm duas coisas:

  \begin{itemize}
    \item a \textbf{integridade} do histórico, já que não dá para trocar o
          passado sem que os identificadores mudem;
    \item o trabalho que dá reescrever o histórico de um repositório
          compartilhado.
  \end{itemize}

  \vspace{0.4cm}
  O campo \textit{pai} é o que dá ao histórico a forma de corrente: cada
  \textit{commit} aponta para o anterior, até o primeiro registro.
\end{frame}

\begin{frame}[fragile]{\texttt{git log}}
  \begin{noterminal}
    \begin{minted}{bash}
git log
git log --oneline
    \end{minted}
  \end{noterminal}

  \begin{minted}{text}
commit 5d09b83191365b932a94e1c90b68b58d5fe12c55
Author: Ana Souza <ana@example.com>
Date:   Sat Aug 8 16:51:53 2026 -0300

    Cria página inicial do catálogo
  \end{minted}

  Ninguém digita os 40 caracteres. Os sete primeiros bastam para identificar o
  \textit{commit}, e é o que o \texttt{\ddash oneline} mostra.
\end{frame}

\begin{frame}[fragile]{Mensagens de commit}
  A mensagem é o único campo que depende de você.

  \begin{minted}{text}
9b1c2d7 update
4a8e330 asdf
c1d0f92 correções
  \end{minted}

  \begin{minted}{text}
9b1c2d7 Corrige alinhamento dos cards em telas estreitas
4a8e330 Adiciona validação de e-mail no formulário de contato
c1d0f92 Substitui imagens placeholder pelas fotos definitivas
  \end{minted}

  O segundo histórico permite achar quando algo entrou sem abrir um único
  arquivo.
\end{frame}

\begin{frame}{Convenções para as entregas da disciplina}
  \begin{itemize}
    \item \textbf{Escreva o que o commit faz}, com verbo no início e no presente
          do indicativo: ``Adiciona'', ``Corrige'', ``Remove'', ``Atualiza''.
    \item \textbf{Uma linha de até cerca de 50 caracteres.} Precisando de mais,
          o \texttt{git commit} sem \texttt{-m} abre um editor: título, linha em
          branco, corpo.
    \item \textbf{Um commit por alteração com sentido próprio.} Nem um
          \textit{commit} por letra digitada, nem um \textit{commit} único no
          fim da aula chamado \texttt{entrega}.
  \end{itemize}
\end{frame}

\begin{frame}[fragile]{Modificado contra não rastreado}
  \begin{noterminal}
    \begin{minted}{bash}
echo '<p>Em construção.</p>' >> index.html
echo 'body { font-family: sans-serif; }' > estilo.css
git status
    \end{minted}
  \end{noterminal}

  \begin{minted}{text}
Changes not staged for commit:
  (use "git add <file>..." to update what will be committed)
  (use "git restore <file>..." to discard changes...)
	modified:   index.html

Untracked files:
	estilo.css
  \end{minted}

  O Git já conhece o \texttt{index.html}, e percebeu que ele mudou. O arquivo
  \texttt{estilo.css} ele nunca viu antes.
\end{frame}

\begin{frame}[fragile]{Preparar um arquivo só}
  \begin{noterminal}
    \begin{minted}{bash}
git add index.html
git commit -m "Adiciona aviso de página em construção"
git status --short
    \end{minted}
  \end{noterminal}

  \begin{minted}{text}
?? estilo.css
  \end{minted}

  \begin{alertblock}{}
    Arquivo que não passou por \texttt{git add} não está no \textit{commit}.
    Essa é a causa mais comum de ``entreguei e o professor disse que faltava
    arquivo''.
  \end{alertblock}
\end{frame}

\begin{frame}[fragile]{\texttt{git status \ddash short}}
  Duas colunas: a primeira é a área de preparação, a segunda é o diretório de
  trabalho. Abaixo, \texttt{\_} representa uma coluna em branco; no terminal
  você verá um espaço.

  \vspace{0.3cm}
  \begin{center}
    \begin{tabular}{ll}
      \toprule
      \textbf{Código} & \textbf{Significado} \\
      \midrule
      \texttt{??} & Arquivo não rastreado \\
      \texttt{\_M} & Modificado, ainda não preparado \\
      \texttt{M\_} & Modificado e preparado \\
      \texttt{A\_} & Arquivo novo, já preparado \\
      \texttt{MM} & Preparado e modificado de novo depois disso \\
      \bottomrule
    \end{tabular}
  \end{center}

  \vspace{0.3cm}
  \texttt{git status} é o comando que você mais vai usar. Execute-o antes de
  cada \texttt{add}, antes de cada \texttt{commit} e depois de cada
  \texttt{push}.
\end{frame}

% ==========================================
\section{Histórico}
% ==========================================

\begin{frame}[fragile]{\texttt{git show}}
  \begin{noterminal}
    \begin{minted}{bash}
git show
    \end{minted}
  \end{noterminal}

  \begin{minted}{text}
commit 7f95d857f84c74ac5325759dcf0ac9de11ed1816
Author: Ana Souza <ana@example.com>
Date:   Sat Aug 8 16:52:03 2026 -0300

    Adiciona aviso de página em construção

diff --git a/index.html b/index.html
index 434dae7..cede13d 100644
--- a/index.html
+++ b/index.html
@@ -1 +1,2 @@
 <h1>Catálogo de Produtos</h1>
+<p>Em construção.</p>
  \end{minted}
\end{frame}

\begin{frame}[fragile]{O formato diff}
  \begin{itemize}
    \item linhas com \texttt{+} foram acrescentadas;
    \item linhas com \texttt{-} foram removidas;
    \item linhas sem marca são contexto, mostradas para você se localizar.
  \end{itemize}

  \vspace{0.4cm}
  A linha \texttt{@@ -1 +1,2 @@} é o cabeçalho do pedaço (\textit{hunk}) que vem
  logo abaixo. Um arquivo com alterações em pontos distantes gera vários
  pedaços, cada um com seu cabeçalho.

  \begin{minted}{text}
@@ -início,quantidade +início,quantidade @@
  \end{minted}

  O \texttt{-} se refere à versão antiga do arquivo, o \texttt{+} à nova.
\end{frame}

\begin{frame}[fragile]{Lendo \texttt{@@ -1 +1,2 @@}}
  \begin{itemize}
    \item \texttt{-1}: no arquivo antigo, o pedaço começa na linha 1 e tem 1
          linha. Quando a quantidade é 1, o Git omite o número.
    \item \texttt{+1,2}: no arquivo novo, o mesmo pedaço começa na linha 1 e
          passou a ter 2 linhas.
  \end{itemize}

  \vspace{0.4cm}
  Em um arquivo maior você veria algo como \texttt{@@ -12,7 +12,9 @@}: sete
  linhas a partir da linha 12 viraram nove linhas a partir da linha 12.

  \vspace{0.4cm}
  Os números localizam a alteração dentro do arquivo. Eles não contam quantas
  linhas mudaram.
\end{frame}

\begin{frame}[fragile]{\texttt{git diff} e \texttt{git diff \ddash staged}}
  \begin{noterminal}
    \begin{minted}{bash}
echo '<p>Segunda linha.</p>' >> index.html
git diff
git add index.html
git diff
git diff --staged
    \end{minted}
  \end{noterminal}

  O segundo \texttt{git diff} não devolve nada. O \texttt{\ddash staged} mostra a
  linha acrescentada.

  \begin{itemize}
    \item \texttt{git diff}: o que eu mexi e ainda não marquei.
    \item \texttt{git diff \ddash staged}: o que vai entrar se eu der
          \texttt{commit} agora.
  \end{itemize}
\end{frame}

\begin{frame}[fragile]{Desfazer antes do commit}
  Tirar da área de preparação, mantendo as alterações no diretório de trabalho:

  \begin{minted}{bash}
git restore --staged index.html
  \end{minted}

  Descartar as alterações e voltar ao estado do último \textit{commit}:

  \begin{minted}{bash}
git restore index.html
  \end{minted}

  \begin{alertblock}{\texttt{git restore} sem \texttt{\ddash staged} apaga o seu
    trabalho e não tem desfazer}
    O conteúdo descartado nunca foi gravado em nenhum \textit{commit}, então não
    existe cópia em lugar nenhum. Confira com \texttt{git diff} antes.
  \end{alertblock}
\end{frame}

\begin{frame}{Desfazer depois do commit}
  Envolve comandos fora do escopo desta disciplina, como o \texttt{git revert} e
  o \texttt{git reset}.

  \vspace{0.5cm}
  Saída suficiente para o nosso caso: faça um \textit{commit} novo que corrija o
  problema.

  \vspace{0.5cm}
  O histórico registra o erro e a correção, o que é honesto e, num diário de
  bordo, é exatamente o que se quer ver.
\end{frame}

\begin{frame}[fragile]{\texttt{.gitignore}}
  \begin{noterminal}
    \begin{minted}{bash}
echo 'minha-senha-secreta' > senha.txt
git status --short
printf 'senha.txt\n*.log\n' > .gitignore
git status --short
git check-ignore -v senha.txt
    \end{minted}
  \end{noterminal}

  \begin{minted}{text}
.gitignore:1:senha.txt	senha.txt
  \end{minted}

  O \texttt{check-ignore} responde qual linha de qual arquivo está escondendo o
  arquivo. Senhas, temporários do editor, saída de compilação e bibliotecas
  baixadas ficam de fora do repositório.
\end{frame}

\begin{frame}[fragile]{Duas ressalvas sobre o \texttt{.gitignore}}
  \begin{itemize}
    \item O próprio \texttt{.gitignore} deve ser versionado, para valer também
          para quem clonar o repositório.
    \item Ele só afeta arquivos \textbf{não rastreados}. Se você já deu
          \texttt{git add} num arquivo, acrescentá-lo ao \texttt{.gitignore}
          depois não o remove do repositório.
  \end{itemize}

  \vspace{0.4cm}
  \begin{noterminal}
    \begin{minted}{bash}
git add .gitignore
git commit -m "Ignora arquivos de senha e de log"
    \end{minted}
  \end{noterminal}
\end{frame}

% ==========================================
\section{Remoto}
% ==========================================

\begin{frame}[fragile]{Repositório remoto}
  Uma cópia do repositório hospedada em outro lugar, que serve de ponto de
  encontro entre as suas máquinas, e entre você e o professor.

  \vspace{0.3cm}
  \begin{center}
    \small
    \begin{tabular}{lll}
      \toprule
      \textbf{Comando} & \textbf{Direção} & \textbf{O que faz} \\
      \midrule
      \texttt{git clone URL} & remoto $\rightarrow$ local & cópia completa, com histórico \\
      \texttt{git push} & local $\rightarrow$ remoto & envia o que só existe aqui \\
      \texttt{git pull} & remoto $\rightarrow$ local & traz o que só existe lá e integra \\
      \bottomrule
    \end{tabular}
  \end{center}

  \vspace{0.3cm}
  Os três transferem \textbf{commits}, e não arquivos soltos. Só vai para o
  servidor aquilo que passou por \texttt{git commit}.

  \vspace{0.2cm}
  \footnotesize \texttt{origin} é o apelido padrão do remoto principal. Não é
  palavra reservada, é convenção.
\end{frame}

\begin{frame}{O endereço do remoto é uma URL, e o transporte é HTTP}
  \begin{center}
    \texttt{https://gitlab.com/ana/catalogo.git}
  \end{center}

  \vspace{0.3cm}
  Esquema, host e caminho, a mesma estrutura vista na aula sobre HTTP. O
  \texttt{git push} num remoto \texttt{https://} faz requisições HTTP, com
  método, cabeçalhos e código de status.

  \vspace{0.4cm}
  Isso explica dois comportamentos que você vai encontrar:

  \begin{itemize}
    \item repositório inexistente ou privado devolve erro de autenticação, e o
          Git para para pedir usuário e senha. A senha não é a da conta: é um
          \textbf{token};
    \item rede do laboratório com a porta bloqueada faz o \texttt{push} falhar
          por tempo esgotado, sem nada de errado no seu repositório.
  \end{itemize}
\end{frame}

\begin{frame}{Forjas}
  Além do Git, um serviço de hospedagem oferece interface web, controle de
  acesso, relatório de problemas e revisão de código. Esse tipo de serviço é
  chamado de \textbf{forja} (\textit{forge}). GitHub, GitLab e Codeberg são
  forjas.

  \vspace{0.4cm}
  \begin{block}{}
    Nenhum desses recursos faz parte do Git. O Git roda na sua máquina; a forja
    é um serviço construído em volta dele. Trocar de forja não altera nenhum dos
    comandos vistos até aqui.
  \end{block}

  \vspace{0.4cm}
  Nesta disciplina usamos o \textbf{GitLab}: o plano gratuito permite
  repositórios privados sem limite de quantidade.
\end{frame}

\begin{frame}{Conta no GitLab}
  \begin{nonavegador}
    \begin{enumerate}
      \item \url{https://gitlab.com/users/sign_up}, com o \textbf{e-mail da
            UFPR} e o seu \textbf{GRR como nome de usuário}, em minúsculas
            (\texttt{grr20249999}). Confirme o e-mail que o GitLab envia;
      \item em \textbf{Edit profile}, preencha o \textbf{nome completo}.
    \end{enumerate}
  \end{nonavegador}

  \vspace{0.4cm}
  O GRR no nome de usuário e o e-mail institucional resolvem a identificação:
  o professor sabe de quem é cada repositório sem perguntar. Nome de usuário do
  tipo \texttt{dark\_coder\_2007} não diz nada a quem corrige.

  \vspace{0.3cm}
  Quem já usa o GitLab, mantenha a conta pessoal como está e crie esta outra
  para a disciplina. Os trabalhos ficam separados dos seus projetos.
\end{frame}

\begin{frame}{A tela \textit{Welcome to GitLab}}
  Depois da confirmação do e-mail, o GitLab pede empresa, grupo e projeto de uma
  vez só, e não deixa pular.

  \vspace{0.2cm}
  \begin{center}
    \small
    \begin{tabular}{ll}
      \toprule
      \textbf{Campo} & \textbf{Resposta} \\
      \midrule
      Company name & \texttt{UFPR} (obrigatório) \\
      Group name & \texttt{ds122-2026-2-turno-grr} \\
      Project name & \texttt{teste} \\
      Role & \texttt{Other} (não há opção de estudante) \\
      Who will be using GitLab? & \texttt{Just me} \\
      Reason for joining & \texttt{I want to learn the basics of Git} \\
      Country or region & \texttt{Brazil} \\
      \bottomrule
    \end{tabular}
  \end{center}

  \vspace{0.2cm}
  \footnotesize O \textbf{Group name} é o que importa: vira o endereço do seu
  grupo e é por ele que as entregas são recolhidas.
\end{frame}

\begin{frame}{Grupo}
  Um \textbf{grupo} reúne vários repositórios sob um mesmo controle de acesso.
  Você usa um só no semestre, e ele já foi criado na tela de boas-vindas.

  \vspace{0.3cm}
  \begin{nonavegador}
    \begin{enumerate}
      \item \textbf{Settings > General}: confira o nome
            \texttt{ds122-2026-2-turno-grr} e a visibilidade \textbf{Private};
      \item \textbf{Manage > Members > Invite members}: procure
            \texttt{alexkutzke} e atribua o papel \textbf{Reporter};
      \item quem não tiver grupo nenhum cria agora, pelo menu \textbf{+},
            \textbf{New group}.
    \end{enumerate}
  \end{nonavegador}

  \vspace{0.2cm}
  \footnotesize O papel \texttt{Reporter} deixa o professor ler e comentar, sem
  poder alterar o seu código.
\end{frame}

\begin{frame}{Fork}
  Uma cópia de um repositório feita \textbf{no servidor}, sob a sua conta ou
  grupo. O original continua onde estava, e você passa a ter um repositório
  próprio, com o mesmo conteúdo e o mesmo histórico, sobre o qual tem permissão
  de escrita.

  \vspace{0.4cm}
  É esse o mecanismo das entregas: o professor publica o repositório com o
  enunciado, cada aluno faz o \textit{fork} para dentro do próprio grupo, e
  resolve ali.

  \vspace{0.4cm}
  \begin{block}{}
    O \textit{fork} acontece no site. O \texttt{clone} acontece na sua máquina.
    A sequência completa é \textit{fork}, \texttt{clone}, \texttt{add},
    \texttt{commit}, \texttt{push}.
  \end{block}
\end{frame}

\begin{frame}{Fork da tarefa de hoje}
  \begin{nonavegador}
    \begin{enumerate}
      \item abra o endereço do repositório da tarefa, publicado na UFPR Virtual;
      \item clique em \textbf{Fork};
      \item em \textbf{Project URL}, escolha \textbf{o seu grupo};
      \item deixe a visibilidade em \textbf{Private} e clique em
            \textbf{Fork project};
      \item copie o endereço \textbf{HTTPS} do seu \textit{fork}.
    \end{enumerate}
  \end{nonavegador}

  \vspace{0.4cm}
  Trabalhando em dupla, apenas um dos dois faz o \textit{fork}, e adiciona o
  colega ao repositório como \texttt{developer}.
\end{frame}

\begin{frame}{A senha da conta não autentica o \texttt{push}}
  Desde abril de 2026 o GitLab exige um segundo fator em todo login feito com
  senha, e o \texttt{git push} não tem como parar para pedir um código.

  \vspace{0.4cm}
  \begin{center}
    \small
    \begin{tabular}{lll}
      \toprule
      \textbf{Onde} & \textbf{Credencial} & \textbf{Endereço do remoto} \\
      \midrule
      Laboratório & token de um dia & \texttt{https://...} \\
      Seu computador & chave SSH & \texttt{git@gitlab.com:...} \\
      \bottomrule
    \end{tabular}
  \end{center}

  \vspace{0.4cm}
  Hoje, aqui, usamos o token. O caminho da chave SSH está no material, em
  \textbf{Autenticação no GitLab}.
\end{frame}

\begin{frame}{Token do dia}
  \begin{nonavegador}
    \begin{enumerate}
      \item em \textbf{janela anônima}, entre no GitLab. Ele envia um código
            para o seu e-mail da UFPR e o pede na tela;
      \item avatar, \textbf{Edit profile}, \textbf{Access > Personal access
            tokens}, \textbf{Generate token};
      \item \textbf{Token name}: \texttt{laboratorio-} e a data de hoje;
      \item \textbf{Expiration date}: amanhã;
      \item \textbf{Scopes}: apenas \texttt{write\_repository};
      \item \textbf{Generate token}, e copie o que aparecer.
    \end{enumerate}
  \end{nonavegador}

  \vspace{0.3cm}
  \begin{alertblock}{}
    O token é exibido \textbf{uma única vez}. Saiu da página, não há como
    recuperá-lo.
  \end{alertblock}
\end{frame}

\begin{frame}[fragile]{Usar o token, e não deixá-lo na máquina}
  \begin{noterminal}
    \begin{minted}{bash}
git config --global credential.helper ""
    \end{minted}
  \end{noterminal}

  O valor vazio desliga o Gerenciador de Credenciais do Windows. Sem ele, o Git
  pergunta a cada \texttt{push} e não grava nada em disco.

  \begin{minted}{text}
Username for 'https://gitlab.com': grr20249999
Password for 'https://grr20249999@gitlab.com':
  \end{minted}

  Usuário é o GRR, senha é o token. Nada aparece enquanto você cola, e colar no
  Git Bash é \texttt{Shift+Insert}.

  \vspace{0.2cm}
  \begin{alertblock}{Antes de sair}
    Revogue o token, apague a pasta clonada e saia da conta no navegador. Toda a
    turma entra nesta máquina como \texttt{Laboratório}.
  \end{alertblock}
\end{frame}

\begin{frame}[fragile]{\texttt{git clone}}
  \begin{noterminal}
    \begin{minted}{bash}
git clone https://gitlab.com/SEU-GRUPO/nome-da-tarefa.git
cd nome-da-tarefa
git remote -v
git log --oneline
git branch --show-current
    \end{minted}
  \end{noterminal}

  \begin{minted}{text}
origin	https://gitlab.com/ana/catalogo.git (fetch)
origin	https://gitlab.com/ana/catalogo.git (push)
  \end{minted}

  O \texttt{origin} já aponta para o seu \textit{fork}, o histórico veio junto e
  o ramo atual é \texttt{main}. Nada disso precisou de \texttt{git init} nem de
  \texttt{git remote add}.
\end{frame}

\begin{frame}[fragile]{Primeiro \texttt{push}}
  \begin{noterminal}
    \begin{minted}{bash}
git config user.name "Seu Nome Completo"
git config user.email "seuemail@example.com"
git add README.md
git commit -m "Registra os integrantes da dupla"
git push
    \end{minted}
  \end{noterminal}

  \vspace{0.3cm}
  \begin{alertblock}{}
    Abra o seu \textit{fork} no navegador e confirme que a alteração está lá.
    Enquanto não aparecer na página do GitLab, a tarefa não foi entregue.
  \end{alertblock}
\end{frame}

\begin{frame}[fragile]{Quando o remoto ainda não conhece o seu ramo}
  Num repositório criado do zero com \texttt{git init}, o primeiro \texttt{push}
  precisa dizer para onde vai:

  \begin{minted}{bash}
git push -u origin main
  \end{minted}

  \begin{minted}{text}
To https://gitlab.com/ana/catalogo.git
 * [new branch]      main -> main
branch 'main' set up to track 'origin/main'.
  \end{minted}

  A opção \texttt{-u} grava a associação entre o seu \texttt{main} local e o
  \texttt{main} do \texttt{origin}. Depois disso, \texttt{git push} e
  \texttt{git pull} sem argumentos já sabem o destino.

  \vspace{0.2cm}
  \footnotesize Quem clonou não precisa disso: o \texttt{clone} já configurou a
  associação.
\end{frame}

\begin{frame}[fragile]{Quando o push é recusado}
  Você deu \texttt{push} de casa na quarta e, na sexta, no laboratório, está com
  uma cópia antiga.

  \begin{minted}{text}
 ! [rejected]        main -> main (fetch first)
error: failed to push some refs to 'https://gitlab.com/...'
hint: Updates were rejected because the remote contains work
hint: that you do not have locally.
  \end{minted}

  A recusa protege o histórico: aceitar o envio apagaria do servidor os
  \textit{commits} que você não tem.

  \begin{minted}{bash}
git pull
git push
  \end{minted}
\end{frame}

\begin{frame}[fragile]{\texttt{divergent branches}}
  Numa instalação nova, o \texttt{git pull} pode parar assim:

  \begin{minted}{text}
hint: You have divergent branches and need to specify how to
hint: reconcile them.
fatal: Need to specify how to reconcile divergent branches.
  \end{minted}

  O Git está perguntando como juntar as duas linhas do tempo. Escolha a
  mesclagem, que é a opção adequada aqui:

  \begin{minted}{bash}
git config --global pull.rebase false
git pull
  \end{minted}

  \vspace{0.2cm}
  \begin{block}{Hábito que evita o problema}
    \texttt{git pull} ao \textbf{começar} a trabalhar, antes de editar qualquer
    arquivo. \texttt{git push} ao terminar.
  \end{block}
\end{frame}

% ==========================================
\section{Conflito}
% ==========================================

\begin{frame}{Quando o Git para e pergunta}
  O \texttt{pull} anterior funcionou porque as duas máquinas mexeram em arquivos
  diferentes.

  \vspace{0.5cm}
  \begin{block}{}
    Quando as duas alteram \textbf{a mesma linha do mesmo arquivo}, o Git não
    tem como decidir, e para.
  \end{block}

  \vspace{0.5cm}
  Vamos provocar o caso de propósito, para você ver a mensagem em condições
  controladas em vez de na véspera da entrega.
\end{frame}

\begin{frame}[fragile]{Provocar o conflito}
  \begin{nonavegador}
    No seu \textit{fork}, edite o \texttt{README.md} pelo próprio site (ícone de
    lápis), altere a primeira linha e confirme a alteração.
  \end{nonavegador}

  \begin{noterminal}
    \begin{minted}{bash}
git commit -am "Altera o título pela máquina local"
git pull
    \end{minted}
  \end{noterminal}

  \begin{minted}{text}
Auto-merging README.md
CONFLICT (content): Merge conflict in README.md
Automatic merge failed; fix conflicts and then commit the result.
  \end{minted}

  Altere \textbf{a mesma linha} que você alterou no site, com outro texto.
\end{frame}

\begin{frame}[fragile]{Onde está o problema}
  \begin{noterminal}
    \begin{minted}{bash}
git status
    \end{minted}
  \end{noterminal}

  \begin{minted}{text}
You have unmerged paths.
  (fix conflicts and run "git commit")
  (use "git merge --abort" to abort the merge)

Unmerged paths:
  (use "git add <file>..." to mark resolution)
	both modified:   README.md
  \end{minted}

  O \texttt{git merge \ddash abort} desfaz a mesclagem inteira e devolve o
  repositório ao estado anterior ao \texttt{pull}. É a saída de emergência.
\end{frame}

\begin{frame}[fragile]{Os marcadores dentro do arquivo}
  Abra o \texttt{README.md} no editor. O Git escreveu as duas versões dentro
  dele:

  \begin{minted}{text}
<<<<<<< HEAD
# Catálogo Verde
=======
# Catálogo de Produtos da Ana
>>>>>>> c34ab3031b9f162880e357de2a5bc27e68d9b6ca
  \end{minted}

  \begin{itemize}
    \item entre \texttt{<<<<<<< HEAD} e \texttt{=======}: a \textbf{sua} versão
          local;
    \item entre \texttt{=======} e \texttt{>>>>>>>}: a versão que veio do
          servidor, identificada pelo \textit{hash} do \textit{commit}.
  \end{itemize}
\end{frame}

\begin{frame}[fragile]{Resolver}
  Editar o arquivo até ele ficar como deve ficar, o que inclui \textbf{apagar as
  três linhas de marcador}. Você pode ficar com uma das versões, com a outra, ou
  escrever uma terceira.

  \begin{noterminal}
    \begin{minted}{bash}
git add README.md
git commit
git push
    \end{minted}
  \end{noterminal}

  O \texttt{git commit} sem \texttt{-m}, aqui, abre o editor com uma mensagem de
  mesclagem já preenchida. Basta salvar e sair.

  \vspace{0.3cm}
  \begin{block}{}
    Conflito não é erro nem defeito do Git. É o Git recusando-se a escolher por
    você entre duas alterações incompatíveis. O erro seria resolvê-lo apagando o
    trabalho do colega sem ler.
  \end{block}
\end{frame}

% ==========================================
\section{Exercícios}
% ==========================================

\begin{frame}{Entrega}
  \begin{block}{Estes exercícios valem nota}
    Compõem o item \textit{Exercícios em sala} da média, e a entrega é até o
    final do encontro de hoje.
  \end{block}

  \vspace{0.4cm}
  A entrega é o próprio \textit{fork} no GitLab, com os \textit{commits} pedidos
  e o \texttt{push} feito. \textbf{Não há link a enviar em lugar nenhum}: o
  professor recolhe os repositórios pelo nome do grupo.

  \vspace{0.4cm}
  Três coisas precisam estar certas: o nome do grupo, o papel de
  \texttt{reporter} do \texttt{alexkutzke} e o \textit{fork} dentro do grupo.
  Grupo fora do padrão não é encontrado, e o que não é encontrado não é
  corrigido.
\end{frame}

\begin{frame}{Como trabalhar}
  A tarefa pode ser feita \textbf{individualmente ou em dupla}. Em dupla, apenas
  um dos dois faz o \textit{fork} no próprio grupo e adiciona o colega ao
  repositório como \texttt{developer}, e a forma de trabalho sugerida é o
  \textbf{Mob Programming}.

  \vspace{0.4cm}
  Registrem no \texttt{README.md} o nome completo e o GRR de cada integrante.

  \vspace{0.4cm}
  \begin{block}{}
    O que está sendo avaliado é o \textbf{histórico} que vocês produzirem, e não
    o conteúdo dos arquivos.
  \end{block}

  \vspace{0.3cm}
  \footnotesize Ao usar IA generativa durante a aula, aplique o pré-prompt
  sugerido no material da disciplina.
\end{frame}

\begin{frame}{Exercícios}
  \begin{block}{Onde está o enunciado}
    O endereço do repositório da tarefa está na atividade da \textbf{UFPR
    Virtual}. O enunciado completo é o \texttt{README.md} desse repositório.
  \end{block}

  \vspace{0.4cm}
  São cinco itens, mais um desafio para quem terminar antes. Todos mexem em
  pouca coisa: editar o \texttt{README.md} e criar alguns arquivos.

  \vspace{0.4cm}
  \begin{alertblock}{}
    Ao terminar, abra o seu \textit{fork} no navegador e confirme que todos os
    \textit{commits} estão lá. \textit{Commit} que não foi enviado não existe
    para quem corrige.
  \end{alertblock}
\end{frame}

% ==========================================
\section*{Resumo}
% ==========================================

\begin{frame}{Resumo}
  \small
  \begin{itemize}
    \item Controle de versão guarda a sequência de estados de um projeto, com
          autoria, data e descrição em cada passo.
    \item O repositório é a pasta \texttt{.git}. O Git funciona sem rede; o
          remoto é opcional.
    \item Um arquivo está no diretório de trabalho, na área de preparação ou no
          repositório. O \texttt{add} avança uma área, o \texttt{commit} avança
          a outra.
    \item Arquivo que não passou por \texttt{git add} não entra no
          \textit{commit}.
    \item O \textit{commit} guarda a fotografia do projeto, o autor, a data, a
          mensagem, o \textit{commit} pai e um identificador calculado a partir
          de tudo isso.
    \item \texttt{git status} responde onde você está; \texttt{git diff} e
          \texttt{git diff \ddash staged} respondem o que mudou, em relação a áreas
          diferentes.
    \item \texttt{clone}, \texttt{push} e \texttt{pull} transferem
          \textit{commits}, nunca arquivos soltos.
    \item Conflito acontece quando duas alterações tocam a mesma linha. O Git
          marca o arquivo e a decisão é sua.
  \end{itemize}
\end{frame}

\begin{frame}{Bibliografia}
  \begin{itemize}
    \item CHACON, S.; STRAUB, B. \textbf{Pro Git}, 2ª ed., cap. 1, 2 e 5. \\
          \url{https://git-scm.com/book/pt-br/v2}
    \item DUDLER, R. \textbf{Guia rápido para o Git}. \\
          \url{http://rogerdudler.github.io/git-guide/index.pt_BR.html}
    \item GitLab. \textbf{Documentação}. \\
          \url{https://docs.gitlab.com/}
  \end{itemize}
\end{frame}

\end{document}
